\section{讨论}
\label{sec:discussion}

\subsection{最重要的发现是一个阴性结果：序列可预测性不可移植}

本文最初的预期是"在严格无泄漏的评估下，序列信号会下降但依然存在"。三个数据集的受控复现实验给出了更尖锐的答案：信号在 DeepCRISPR 上弱（中位 $R^{2}$ $0.067$--$0.120$）、在 Hiranniramol 上强（$0.345$--$0.489$）、在 Labuhn 上\textbf{完全不存在}（全部为负）。而且这不是划分造成的假象——一个与官方划分完全无关的全数据交叉验证给出同向结论，Labuhn 的 CV $R^{2}$（$-0.144$）仅比它自己的标签置换零分布（$-0.287\pm0.050$）高 $0.14$。

这一结果的实践含义是直接的：\textbf{"某模型在 sgRNA 效率预测上达到 $R^{2}=0.7$"这类陈述，不能作为关于 sgRNA 设计的一般性知识被引用}，除非同时说明它在哪些独立数据集上被复现、以及复现失败的比例。这也解释了为什么独立基准测试反复观察到"报告性能在严格评估下明显下降"\citep{Zhang2023BriefBioinform,Konstantakos2022NAR}：问题可能不只是评估不够严格，而是\textbf{可预测性本身强烈依赖具体的数据集}。

本文因此把"复现"从一句声明变成了一个需要逐数据集报告的结果。我们刻意保留了复现失败的数据集，而不是只报告成功的那个——后者会让结论看起来强得多，但也会让读者无法判断该结论的适用边界。

\subsection{归因需要外部系统才能成为可检验的陈述}

多模型归因在文献中常被当作"解释"的终点：如果线性系数、SHAP、IG 与注意力都指向同一个位置，似乎就可以说该位置重要。本文的做法揭示了两点。

\textbf{第一，多模型一致本身不能排除共同混杂。} 本文五类模型共享同一份数据、同一份标签与同一套预处理；它们一致指向第 18 位，原则上也可能是因为它们共同依赖了数据中某个与第 18 位相关的结构（例如该位点的碱基分布与 GC 含量、与 locus 类型、与细胞系组成的耦合）。仅凭 E3 无法区分"第 18 位重要"与"第 18 位是某个混杂的代理"。

\textbf{第二，换一个系统可以明显提高证据强度。} 当我们用 CRISPRon——一个由他人训练、使用 30\,nt 窗口与 CRISPRoff 能量项、输出 0--100 复合评分的模型——去预测同一个最小扰动时，两个系统在 8 个预登记实例中的 7 个上方向一致，且变化幅度正相关（$r=0.629$）。这一步不能证明第 18 位在生物学上重要（那需要 E5），但它确实把结论从"我们的模型这样做"推进到"两个独立构造的模型系统对同一扰动给出同向响应"。这是 E3 无法单独达到的强度。

我们特意保留了两处"不好看"的结果，因为它们是这一结论可信度的条件：唯一分歧样本 WT08 被保留在统计中；饱和突变虽然把第 18 位排在第 1，但第 2 名（第 17 位，7.68）与第 3 名（第 20 位，7.18）并不遥远，说明第 18 位的突出是\textbf{量级上的}而非\textbf{类别上的}。

\subsection{"未检出"是结果，不是失败}

本文报告了三类阴性或未检出结果：环境通道无增量预测证据（四因子均 Inconclusive）、跨细胞系迁移失败（LOCO 中位 $R^{2}$ 与 0 不可区分）、Labuhn 上无序列信号。它们各自有不同性质，不应被笼统地归为"模型不行"。

\textbf{环境通道的结果有明确的可检测性上限。} 四个因子的 $|\overline{\Delta R^{2}}|$ 为 $3.2\times10^{-4}$--$1.5\times10^{-3}$，比预设门 $0.01$ 低 1--2 个数量级，跨模型区间全部跨 0，析因 ANOVA $p=0.975$--$0.9996$。三者一致。但这一"未检出"受两个具体限制约束：环境通道只有 4 条二值轨道，不足以代表完整 cell state；若干细胞系--通道组合近乎常数，直接压缩了可检测效应上限。因此正确的陈述是"在通用多细胞系设定与该编码分辨率下未检出"，而不是"表观遗传信息与编辑效率无关"——已有研究在特定细胞类型中报告了表观特征带来的增益\citep{Ito2024NAR,Luboshitz2026TCBB}，而另一些研究报告加入 DNase 轨道后并未提升性能\citep{Wang2019NatCommun}。本文的结果与后者一致。

\textbf{跨细胞系迁移的失败夹杂了训练集缩减。} LOCO 的训练池因剔除留出系同源身份类而缩减至 1\,453--12\,313 条，"分布漂移"与"样本减少"两种因素在本文数据中不可分离（\S\ref{sec:limitations}）。因此本文只报告"与均值基线不可区分"，不报告"跨细胞系必然失败"。

\subsection{把"能说什么"写进方法：六层证据的实际作用}

六层证据分类（表~\ref{tab:evidence-layers}）不是装饰性的，它在本文中至少三次改变了结论的写法。

\textbf{第一次：防止 E3 冒充统计证据。} 归因幅值在本文中被明确排除在证据等级判定之外，SNR 也明确为辅助标注而非检验统计量（\S\ref{sec:methods}）。因此"第 18 位归因最高"这句话只能写成描述性汇总，不能写成"显著"。

\textbf{第二次：防止 E4 冒充 E5。} 外部模型验证得到 7/8 的一致性，这是一个相当整齐的结果。如果没有证据分层，很容易被写成"第 18 位替换降低编辑效率"。本文统一表述为"两个模型系统对同一扰动的预测变化方向一致"，并明确指出这仍是模型行为。WT08 的分歧也因此在文中被如实保留，而不是作为噪声处理。

\textbf{第三次：防止 E2 与 E3/E4 的"三方一致"被升级为因果。} 三条线索方向一致确实提高了候选的可信度，但 E2 自身就揭示了反例（HEK293T 上方向反转）。三条线索不能相互替代：E2 是观测关联，E3 是模型依赖，E4 是模型扰动。它们在同一位置的汇合说明该位置值得实验检验，仅此而已。

\subsection{平台定位：不是更高精度，而是更清楚的分界}

本文\textbf{不}主张平台在预测精度上优于已有方法。DeepCRISPR 主数据集上的中位 $R^{2}$ 仅 $0.067$--$0.120$，跨细胞系迁移失败，环境通道无增益；这些数字都低于文献中常见的报告值，而本文认为这种"低"更接近严格评估下的真实水平。平台的贡献在于把"哪些结论在什么证据强度上成立"变成流程的输出：数据质量审计给出 E1 的可信度前提，跨数据集复现给出 E1 的适用边界，归因与外部验证给出 E3/E4 的收敛点，观测性关联给出 E2 的独立对照，候选表给出 E5 的实验设计，而证据分层确保这五者不会被混为一谈。

这也带来一个方法学上的连带结果：本文在整改自身历史批次的泄漏时，发现并修正了一个更隐蔽的问题——同一批 run 同时产出测试集与验证集指标，而审计脚本依赖文件枚举顺序取值，导致数值发散计数在两个口径间漂移（33 与 20）。该缺陷不影响任何模型训练，但会影响"结果是否可信"这类元数据的准确性。我们把它写进论文而不是静默修掉，因为它正属于本文主张要审计的那一类问题：\textbf{报告数字的口径必须被显式固定并被代码强制。}
